\chapter{Теоретические сведения}

Данная глава вводит основные понятия, необходимые для понимания методики и результатов работы. В ней рассматриваются внутренние представления нейронных сетей, частотный анализ изображений, свойства спектра, архитектуры ResNet, Vision Transformer и Swin Transformer, а также исследования, связывающие частотные характеристики с обучением и интерпретируемостью моделей компьютерного зрения.

\section{Внутренние представления и частотная интерпретация}

\subsection{Внутренние представления моделей компьютерного зрения}

Глубокую нейронную сеть можно рассматривать как последовательность преобразований входных данных. В случае компьютерного зрения на вход модели подаётся изображение, а на выходе получается предсказание: например, класс объекта. Однако между входом и выходом находится большое число промежуточных слоёв. Каждый слой преобразует данные в новую форму. Эти промежуточные результаты называются внутренними представлениями, признаками или активациями.

Интуитивно внутреннее представление — это то, как сама модель «видит» изображение на некотором этапе обработки. На ранних слоях признаки обычно ближе к исходному изображению: они могут реагировать на границы, контуры, простые текстуры и локальные цветовые переходы. На более глубоких слоях признаки становятся более абстрактными: отдельные нейроны или каналы могут реагировать уже не на конкретный пиксельный рисунок, а на части объектов, повторяющиеся структуры или семантические категории. Такая иерархия признаков является одной из причин эффективности глубоких моделей.

Для свёрточных сетей внутренние представления обычно имеют вид карт признаков. Карта признаков — это двумерная сетка значений, похожая на изображение, но вместо исходных пикселей она содержит отклики фильтров. Если у изображения есть три канала цвета, то у карты признаков глубокого слоя могут быть десятки, сотни или тысячи каналов. Каждый канал можно понимать как отдельный тип признака, который модель научилась выделять.

Для трансформеров представления имеют другую форму. Визуальный трансформер разбивает изображение на патчи, то есть небольшие непересекающиеся фрагменты. Каждый патч превращается в вектор — токен. После этого модель работает не с изображением как с двумерной сеткой пикселей, а с последовательностью токенов. Поэтому анализ трансформеров требует дополнительной интерпретации: чтобы сравнить их с CNN, токены можно вернуть в двумерную решётку, соответствующую расположению патчей в исходном изображении.

Внутренние представления важны, потому что именно они определяют поведение модели. Две архитектуры могут иметь похожую точность классификации, но при этом строить совершенно разные признаки. Одна модель может сильнее опираться на текстуры, другая — на форму объекта, третья — на глобальный контекст. Поэтому анализ внутренних представлений позволяет лучше понять, почему модель работает именно так, как она работает.

Классическая работа \cite{yosinski2014transferable} показывает, что признаки разных уровней обладают разной степенью переносимости между задачами. Ранние признаки часто оказываются более универсальными, тогда как поздние сильнее связаны с конкретной задачей и датасетом. В данной работе этот взгляд дополняется частотной интерпретацией: нас интересует не только то, насколько признаки переносимы, но и то, как по глубине модели меняется их спектральный состав.

\subsection{Изображение как пространственный сигнал}

Для понимания спектрального анализа важно рассматривать изображение не только как набор пикселей, но и как пространственный сигнал. Пространственный сигнал — это функция, которая задана на координатной сетке. В каждой точке изображения есть соответствующее значение цвета или яркости, в зависимости от используемой цветовой модели. 
При движении по координатной сетке может быть плавное и медленное изменение значения от пикселя к пикселю, а может резко изменяться - первое поведение соответствует низким пространственным частотам, второе - высоким пространственным частотам.

Плавный фон, большое однотонное пятно, объект крупной формы - все это пример низкочастотной структуры. Резкие границы, тонкие линии, мелкие детали и текстуры и шум - высокочастотная структура. Реальное изображении содержит в себе одновременно компоненты разных частот. Например, фотография собаки на заднем сиденье содержит низкочастотную информацию о силуэте и общей форме животного и фона, высокочастотную информацию о шерсти, краях и мелких деталях.

Для решения задач компьютерного зрения важно это учитывать. Для одних задач достаточно низких частот: например, чтобы грубо отличить автомобиль от самолёта, часто хватает общей формы. Для других задач высокие частоты критичны: в медицинских снимках, спутниковых изображениях или дефектоскопии мелкие детали могут нести ключевую информацию. Поэтому способность модели сохранять или подавлять разные частотные компоненты влияет на её применимость к конкретным задачам.

Частотный анализ позволяет перейти от вопроса «какие значения имеют пиксели или активации?» к вопросу «какие масштабы изменений представлены в сигнале?». В этой работе такой подход применяется не только к исходным изображениям, но и к промежуточным представлениям нейронных сетей.

\subsection{Спектр, частотный радиус и радиальное усреднение}

Спектр — это описание сигнала через набор частотных компонент. В пространственной области мы смотрим, как значения сигнала расположены по координатам. В частотной области мы смотрим, из каких колебаний разного масштаба этот сигнал состоит. Можно сказать, что спектр показывает, сколько в сигнале «медленных» и «быстрых» изменений.

Для изображения низкие частоты отвечают за крупные и плавные структуры. Средние частоты часто связаны с частями объектов, локальными формами и выраженными, но не слишком мелкими деталями. Высокие частоты отвечают за резкие изменения, мелкие текстуры, тонкие контуры и шум. Поэтому спектр изображения можно использовать как способ описать, насколько изображение или карта признаков являются гладкими, детальными или шумными.

В спектральном анализе обычно различают амплитуду, фазу и энергию. Амплитуда показывает силу частотной компоненты. Фаза содержит информацию о расположении соответствующей структуры в пространстве. Энергия показывает вклад частоты в общий сигнал. В данной работе, поскольку задача состоит в том, чтобы понять, какие частотные диапазоны сохраняются в представлениях разных слоёв, основной интерес представляет распределение энергии по частотам.

Также важной составляющей любого изображения является нулевая частота или, так называемая, DC-компонента. нулевая частота -  это постоянная составляющая сигнала, то есть его средний уровень. DC-компонента может оказаться очень большой и скрыть более содержательные различия между частотами, если среднее значение карты признаков имеет большое положительное значение. Поэтому в практическом анализе часто удаляют пространственное среднее перед вычислением спектра. Процедура этого шага подробно описана в главе с методикой работы.

По одному только спектру нельзя решить, полезна ли конкретная частота для классификации. Высокочастотная область может соответствовать границам и мелким деталям, но там же могут находиться шумовые или нестабильные компоненты. Низкие частоты часто отражают форму и крупные структуры, однако для некоторых задач такого описания слишком мало. Поэтому спектральный анализ в данной работе используется как способ описать представления, а не как самостоятельная мера качества модели

Изучение двухмерных карт признаков или изображений предполагает использование пространственного спектра с аналогичной размерностью. В двумерном спектре фиксируется амплитуда колебаний и направление: вертикальное, горизонтальное, диагональное и другие. Например, вертикальные и горизонтальные переходы могут характеризоваться сопоставимым пространственным масштабом, но различаться по ориентации в частотной плоскости. 

Направление не имеет значения, поэтому для анализа важен лишь частотный масштаб, из-за чего исходный двумерный спектр преобразуется в радиальную одномерную кривую. Для этого частоты с одинаковым расстоянием от центра спектра усредняются. Получается зависимость энергии от частотного радиуса. Малый радиус соответствует низким частотам, большой радиус — высоким.

Радиальный спектр удобен для сравнения разных моделей и слоёв. Вместо сложной двумерной карты мы получаем одну кривую, по форме которой можно судить о частотном профиле представления. Если кривая быстро убывает после пика в низких частотах, значит представление сильно сконцентрировано в низкочастотной области. Если хвост кривой тянется дальше, значит в представлении сохраняется вклад средних или высоких частот.

В результатах работы используется нормализованный частотный радиус. Это нужно для того, чтобы сравнивать слои с разным пространственным разрешением. Например, ранний слой CNN может иметь большую карту признаков, а поздний — маленькую. У таких карт различается максимальная представимая частота. Нормализация частотной оси позволяет привести кривые к сопоставимому виду. В экспериментальной части дополнительно используется приведение частотной шкалы к масштабу первого слоя, что позволяет сравнивать спектры разных блоков на общей оси.

\subsection{Интерпретация низких и высоких частот}

Спектральный анализ даёт удобный способ описывать внутренние представления модели, но трактовать его результаты нужно аккуратно. Прежде всего, спектральная энергия не является информацией в строгом теоретико-информационном смысле. Она показывает, как распределены активации по частотным диапазонам, но сама по себе не отвечает на вопрос, насколько эти признаки полезны для распознавания объекта.

То же относится и к разделению частот на низкие и высокие. Высокочастотная область может содержать действительно важные детали: границы, мелкие текстуры, локальные различия между объектами. Но в эту же область часто попадают шумовые и нестабильные компоненты. С низкими частотами ситуация похожая. Они могут хорошо передавать общую форму и крупную структуру объекта, однако для задач, где решающими оказываются небольшие отличия, такого описания может быть недостаточно. Поэтому в работе частотный профиль рассматривается не как прямой показатель качества, а как дополнительная характеристика представлений.

Отдельная осторожность нужна при сравнении CNN и трансформеров. Карта признаков свёрточной сети и решётка токенов ViT устроены по-разному. В CNN пространственные карты формируются последовательностью свёрточных фильтров. В ViT каждый токен соответствует целому патчу и уже содержит агрегированную информацию о фрагменте изображения. Возврат токенов к двумерной решётке позволяет применить общий спектральный аппарат, но не делает представления этих архитектур полностью одинаковыми.

Есть и ещё одно ограничение. В данной работе анализируются уже обученные модели, то есть фиксируется состояние представлений после завершения обучения. По таким данным можно сравнить финальную спектральную структуру слоёв, но нельзя проследить, как она возникала по эпохам. Для этого потребовался бы отдельный эксперимент: сохранять промежуточные состояния модели во время обучения и анализировать, как меняются спектральные характеристики на разных этапах.

В обычной обработке изображений низкие частоты часто связывают с общей формой и плавными областями, а высокие — с деталями и текстурами. В нейронных сетях эта интерпретация сохраняется, но становится более сложной. Внутренние представления — это уже не исходные пиксели, а признаки, сформированные моделью. Поэтому низкие частоты в карте признаков означают не просто плавность исходного изображения, а пространственную гладкость некоторого признака. Высокие частоты означают, что признак быстро меняется по пространственным координатам.

Например, если некоторый канал карты признаков активируется на всей области объекта, его спектр будет более низкочастотным. Если канал активируется только на узких границах или мелких деталях, в нём будет больше высокочастотных компонент. Если глубокий слой теряет высокочастотные компоненты, это может означать, что модель переходит от локальных деталей к более глобальным и устойчивым признакам. Но это также может означать потерю информации, важной для задач, чувствительных к мелким структурам.

Именно поэтому в данной работе используется термин «спектральная энергия представления». Он не означает информацию в строгом смысле теории информации Шеннона. Речь идёт о том, как энергия активаций распределена по пространственным частотам. Такой подход позволяет исследовать частотное разнообразие признаков, но не утверждает напрямую, что одна частота «важнее» другой для классификации.

Изменение частотного состава по глубине сети можно рассматривать как один из признаков того, как модель переходит от локальной информации к более обобщённым признакам.

В литературе показано, что частотные компоненты изображения могут быть связаны с поведением моделей компьютерного зрения. Например, в работе \cite{yin2019fourier} частотная область используется для анализа устойчивости моделей к различным искажениям, а в \cite{wang2020highfrequency} обсуждается роль высокочастотных компонент в обобщающей способности CNN. Эти результаты важны для данной работы не как непосредственный объект проверки, а как дополнительное обоснование того, что частотный состав представлений может отражать существенные свойства модели.

Для визуальных трансформеров вопрос частотной интерпретации также остаётся актуальным. Работы \cite{naseer2021intriguing,paul2022vision,alijani2024vision} показывают, что ViT могут отличаться от CNN по характеру устойчивости и внутренней организации признаков. Это согласуется с общей мотивацией настоящего исследования: сравнивать архитектуры следует не только по итоговой точности, но и по тому, как они преобразуют признаки внутри модели.


\section{Архитектуры компьютерного зрения}

\subsection{Свёрточные нейронные сети и ResNet}

Свёрточная нейронная сеть — это архитектура, специально приспособленная для обработки данных с пространственной структурой. В изображении соседние пиксели обычно связаны друг с другом, и CNN использует это свойство через локальные фильтры. Фильтр смотрит на небольшой участок изображения или карты признаков и вычисляет отклик. Один и тот же фильтр применяется во всех пространственных позициях, поэтому сеть может находить один и тот же признак в разных местах изображения.

Для понимания CNN важны несколько терминов. Канал — это отдельная карта признаков, соответствующая одному типу выделенного паттерна. Ядро свёртки — это небольшой набор весов, который перемещается по изображению. Stride — шаг, с которым фильтр проходит по пространственным координатам. Если шаг больше единицы, пространственное разрешение уменьшается. Pooling — операция агрегирования соседних значений, также часто уменьшающая разрешение. Рецептивное поле — область исходного изображения, которая влияет на конкретное значение в карте признаков. Чем глубже слой, тем больше обычно его рецептивное поле.

С точки зрения спектрального анализа CNN можно понимать как систему фильтров. Ранние фильтры часто реагируют на простые локальные структуры: границы, углы, текстуры. Такие признаки могут содержать заметный вклад средних и высоких частот. Более глубокие слои объединяют локальные признаки в более крупные структуры, и их карты признаков становятся более низкочастотными. Дополнительно операции уменьшения разрешения физически ограничивают набор частот, которые могут быть представлены на карте признаков.

Именно поэтому для CNN естественно ожидать спектральное сжатие по глубине. Под спектральным сжатием в данной работе понимается смещение энергии представлений к низким частотам и уменьшение вклада более высоких частотных диапазонов. В результатах это проявляется через форму спектральных кривых и численные метрики.

Семейство ResNet относится к свёрточным архитектурам, в которых ключевую роль играют остаточные блоки \cite{he2016deep}. До появления ResNet увеличение числа слоёв в CNN не всегда давало выигрыш: глубокая модель могла хуже оптимизироваться и показывать результат ниже, чем более короткая сеть. Эта проблема не сводилась только к переобучению. Существенную роль играла сложность передачи полезного градиента через длинную цепочку преобразований.

Остаточная связь меняет постановку задачи для блока. Вместо того чтобы каждый фрагмент сети заново строил полное преобразование входа, он обучает добавочное изменение к уже имеющемуся представлению. За счёт такого обходного пути сигнал может проходить через сеть устойчивее, а обучение глубоких моделей становится практически выполнимым. Именно поэтому ResNet стала одной из базовых архитектур для анализа глубины в компьютерном зрении.

В данной работе ResNet используется как представитель классической иерархической CNN. У этого выбора есть методическое преимущество: внутри одного семейства доступны модели разной глубины — ResNet18, ResNet34, ResNet50 и ResNet101. Это позволяет отделить влияние общей архитектурной идеи от влияния масштаба модели. Если сравнивать разные семейства сразу, трудно понять, что именно изменило спектральный профиль: свёртки, внимание, глубина или способ уменьшения разрешения. Сравнение вариантов ResNet частично снимает эту неоднозначность.

Структура ResNet также удобна для послойного анализа. Обычно в ней выделяют крупные группы блоков: layer1, layer2, layer3 и layer4. Это не отдельные операции, а целые стадии обработки. При переходе между ними меняются пространственный размер карт признаков, число каналов и эффективное рецептивное поле. Поэтому такие уровни хорошо подходят для исследования крупной динамики представлений: от ранних карт, ещё сохраняющих пространственную детализацию, к более глубоким признакам, где локальные изменения уже сильнее агрегированы.

\subsection{Transformer и механизм самовнимания}

Transformer возник как архитектура для обработки последовательностей, в которой основной вычислительный механизм строится вокруг внимания \cite{vaswani2017attention}. В CNN связь между элементами задаётся локальным фильтром: каждый фильтр смотрит на небольшой участок карты признаков. В Transformer другая логика. Элемент последовательности может учитывать другие элементы не через заранее заданное соседство, а через вычисляемые моделью веса взаимодействия.

Механизм self-attention можно описать как адаптивное смешивание токенов. Для каждого токена модель оценивает, какие элементы последовательности важны для его обновления. Если связь между токенами существенна, соответствующий вклад становится больше; если связь слабая, вклад уменьшается. Multi-head attention расширяет эту схему: несколько голов внимания работают параллельно и могут кодировать разные типы отношений между токенами.

Обычный трансформерный блок не ограничивается attention. В него входят нормализации, полносвязный блок, часто называемый MLP или feed-forward network, а также остаточные соединения. Attention отвечает за обмен информацией между токенами. MLP преобразует признаки каждого токена. Нормализация стабилизирует масштабы активаций. Остаточные связи сохраняют часть прежнего представления и помогают обучению глубоких стеков блоков. Поэтому поведение Transformer нельзя объяснять только матрицами внимания: итоговый выход блока возникает из совместной работы нескольких механизмов.

Для частотного анализа эта деталь принципиальна. Смешивание токенов действительно может вести себя как сглаживающая операция, поэтому в литературе self-attention связывают с низкочастотной фильтрацией \cite{park2022how}. Но в реальной архитектуре рядом с attention находятся MLP, LayerNorm и residual-соединения. Они способны менять спектральный профиль иначе, чем простой оператор усреднения. По этой причине вопрос о том, будет ли визуальный трансформер монотонно сдвигать энергию к низким частотам, нельзя решить только из описания блока. Его нужно проверять на промежуточных представлениях модели.

\subsection{Vision Transformer}

Vision Transformer переносит трансформерный подход на изображения \cite{dosovitskiy2021image}. Изображение сначала разбивается на фиксированные непересекающиеся фрагменты — патчи. Например, при размере входа $224 \times 224$ и размере патча $16 \times 16$ получается решётка из 196 патчей. Каждый такой фрагмент разворачивается в вектор и проходит через линейную проекцию. После этого патч рассматривается уже не как часть изображения, а как токен в последовательности.

У такой последовательности сама по себе нет знания о двумерном расположении элементов. Чтобы модель различала, где находился каждый патч, к токенам добавляют позиционные эмбеддинги. В классификационных вариантах ViT обычно присутствует также CLS-токен. Он не соответствует конкретному участку изображения; его задача — собрать информацию, необходимую для финального предсказания класса.

Различие между ViT и CNN затрагивает сам способ обработки изображения. В CNN локальность заложена конструктивно: фильтр работает с соседними пикселями или соседними элементами карты признаков. В ViT такого жёсткого ограничения нет. Через механизм внимания токен может получать информацию от других патчей, в том числе удалённых. Это позволяет модели раньше учитывать глобальный контекст, но одновременно ослабляет те индуктивные смещения, которые у CNN связаны с локальными текстурами и постепенной пространственной агрегацией.

Сравнение внутренних представлений ViT и CNN отдельно изучалось в работе \cite{raghu2021do}. Авторы показывают, что визуальные трансформеры иначе распределяют информацию по глубине: они раньше вовлекают глобальные связи и иначе используют остаточные соединения. Для данной работы это означает, что спектральная динамика ViT не обязана повторять поведение ResNet. Даже если attention обладает сглаживающими свойствами, полный стек ViT может перераспределять частотную энергию сложнее.

В экспериментальной части выходы ViT рассматриваются как решётка патчевых токенов. Для этого CLS-токен исключается, а оставшаяся последовательность возвращается к двумерному порядку патчей. Такой приём нужен, чтобы применить к ViT тот же частотный аппарат, что и к CNN. При этом интерпретация требует осторожности: токен ViT соответствует целому патчу, а не отдельной пространственной позиции пиксельного уровня. Поэтому решётка токенов является удобным приближением для анализа, но не полной копией карты признаков свёрточной сети.

\subsection{Swin Transformer}

Swin Transformer был предложен как попытка адаптировать трансформерную архитектуру к особенностям изображений более явно, чем это делает классический ViT \cite{liu2021swin}. Главная проблема классического ViT состоит в том, что глобальное внимание между всеми токенами становится вычислительно дорогим при больших разрешениях. Кроме того, изображения имеют естественную иерархию объектов разных масштабов, а классический ViT не строит такую иерархию так явно, как CNN.

Swin решает эти проблемы через два ключевых механизма. Первый — оконное внимание. Вместо того чтобы каждый токен взаимодействовал со всеми токенами изображения, внимание вычисляется внутри локальных окон. Это снижает вычислительную сложность и вводит локальность, похожую на локальность свёрток. Второй механизм — сдвиг окон, или shifted windows. На соседних блоках окна немного смещаются, благодаря чему информация может передаваться между областями, которые раньше находились в разных окнах.

Третий важный элемент — patch merging. Это операция объединения соседних патчей, уменьшающая пространственное разрешение и увеличивающая размерность признаков. По смыслу она похожа на понижение разрешения в CNN. Благодаря patch merging Swin строит иерархию признаков: ранние стадии имеют более высокое разрешение, поздние — более низкое, но более богатое каналами представление.

С точки зрения спектрального анализа Swin особенно интересен, потому что сочетает свойства CNN и ViT. С одной стороны, он использует механизм внимания и токенные представления. С другой стороны, он имеет локальность, стадии и уменьшение пространственного разрешения. Поэтому можно ожидать, что его спектральная динамика будет промежуточной: не полностью такой же, как у ResNet, но и не такой же, как у классического ViT.

Описанные различия между ResNet, ViT и Swin определяют выбор экспериментальной схемы в следующей главе. Для ResNet анализируются карты признаков, для ViT — решётки патчевых токенов, для Swin — токенные представления с учётом стадий и уменьшения пространственного разрешения.

\section{Спектральные свойства нейронных сетей}

\subsection{Спектральное смещение и принцип частот}

Одно из важных направлений исследований в теории глубокого обучения связано с понятием spectral bias, или спектрального смещения. Под этим понимается склонность нейронных сетей легче и быстрее обучать низкочастотные зависимости, чем высокочастотные. Иначе говоря, модель сначала хорошо приближает более гладкие и простые компоненты функции, а более резкие и быстро меняющиеся зависимости осваивает позднее или хуже.

Работа \cite{rahaman2019spectral} является одной из ключевых в этом направлении. Авторы показали, что глубокие сети с ReLU-активациями обладают смещением в сторону низких частот. Это не означает, что нейронная сеть вообще не может выучить высокочастотную функцию. Скорее это означает, что низкочастотные закономерности оказываются для неё более естественными с точки зрения процесса обучения.

Позднее эта идея развивалась в рамках frequency principle, или принципа частот, где изучается более общий принцип освоения частотных компонент в процессе обучения \cite{xu2025overview}. В работе \cite{kiessling2022computable} предложены способы вычислимого определения спектрального смещения, то есть переход от качественного описания к более формальному измерению. Работа \cite{tancik2020fourier} показывает, что использование Фурье-признаки может помочь сетям лучше представлять высокочастотные функции, что дополнительно подчёркивает значимость частотной области.

Для данной диссертации эти работы важны как теоретическая основа. Авторы в своих работах подтверждают, что спектральные свойства глубоких моделей закономерна и не является артефактом. Напротив, они отражают как модели обучаются, как формируют внутренние представления и какие признаки и структуры усваивают в первую очередь.


\subsection{Attention как фильтрация признаков}

Механизм внимания можно обсуждать в частотных терминах потому, что он меняет пространственное распределение признаков. Когда блок смешивает токены, близкие или удалённые, он фактически перераспределяет значения по решётке представления. Такое смешивание иногда ведёт к сглаживанию, а сглаживание в частотной области обычно связано с ослаблением высоких частот. По этой причине self-attention в ряде работ рассматривается как оператор с низкочастотными свойствами.

"В работе \cite{park2022how} показано, что self-attention может вести себя как низкочастотный оператор. Это объясняет часть эффектов, наблюдаемых в ViT, в том числе склонность к сглаживанию при увеличении глубины. Но переносить этот вывод на весь блок ViT напрямую нельзя. Помимо attention в нём есть многослойный перцептрон, нормализация слоя, остаточные соединения и позиционные эмбеддинги. Их совместное действие может дать немонотонный спектральный профиль, что как раз проверяется в экспериментальной части.

Работа \cite{bai2022improving} дополнительно показывает, что высокочастотные компоненты могут быть важны для улучшения ViT. Это означает, что частотный анализ может не только объяснять поведение моделей, но и подсказывать направления для изменения архитектур или обучения.

\subsection{Частотно-ориентированные методы в глубоком обучении}

Частотная область используется не только для анализа, но и для построения новых методов. В работе \cite{xu2020learning} рассматривается обучение в частотной области. Идея таких подходов состоит в том, что некоторые закономерности удобнее представлять не в пространственной форме, а через частотные компоненты. Работа \cite{li2021frequency} демонстрирует использование слияние частотного домена в CNN для задач диагностики и адаптации домена. Это показывает, что частотные признаки могут быть полезны не только в классическом компьютерном зрении, но и в прикладных инженерных задачах.

В более широком контексте глубокое обучение всё чаще применяется к спектральному анализу, моделированию и восстановлению сигналов \cite{liu2024deep}. Для трансформеров также появляются частотно-ориентированные архитектуры и модификации. Обзоры \cite{aslan2026frequency,li2025review} рассматривают модели, использующие преобразования Фурье, вейвлеты, дискретное косинусное преобразование и другие инструменты частотной области в задачах обработки изображений.

Для данной работы эти исследования важны как подтверждение того, что частотная область является не только вспомогательной математической техникой, но и активным направлением развития архитектур глубокого обучения. Анализ спектральных характеристик представлений может быть первым шагом к более осознанному проектированию моделей, которые лучше сохраняют или подавляют нужные частотные диапазоны.

\subsection{Датасет ImageNet и классификация}

ImageNet является одним из ключевых датасетов в истории компьютерного зрения \cite{deng2009imagenet}. Он содержит большое количество изображений, распределённых по множеству классов, и долгое время использовался как основной ориентир для сравнения архитектур классификации. Многие модели CNN и трансформеров предобучаются или оцениваются именно на ImageNet, поэтому этот датасет удобен для сопоставления архитектур.

В данной работе используется подвыборка ImageNet 10K \cite{priyerana_imagenet10k}. Это позволяет проводить анализ на реальных изображениях из крупномасштабного домена, сохраняя вычислительную выполнимость экспериментов. Важно, что цель работы состоит не в обучении моделей с нуля, а в анализе уже сформированных представлений. Поэтому использование валидационных изображений ImageNet позволяет исследовать поведение архитектур на разнообразных естественных изображениях.

Крупномасштабная классификация важна ещё и потому, что на таких данных модели вынуждены учитывать как глобальные формы, так и локальные детали. Одни классы могут различаться по общей форме объекта, другие — по текстуре или мелким признакам. Поэтому спектральный анализ представлений на ImageNet даёт более содержательную картину, чем анализ на искусственных или слишком простых данных.


В задачах классификации качество модели обычно измеряется точностью. Top-1 accuracy показывает долю изображений, для которых наиболее вероятный класс, предсказанный моделью, совпал с правильным классом. Эта метрика удобна и понятна, но она не объясняет, какие признаки использует модель.

Спектральные характеристики дают другую информацию. Они показывают, как распределена энергия представлений по частотам. Модель с низким спектральным центроидом может сильнее опираться на сглаженные и глобальные признаки. Модель с более широким спектром может сохранять больше локальных деталей. Но из этого нельзя напрямую сделать вывод, что одна модель лучше другой. Всё зависит от того, какие признаки нужны для конкретной задачи.

Поэтому в данной работе accuracy и спектральные метрики рассматриваются как взаимодополняющие характеристики. Accuracy отвечает на вопрос, насколько хорошо модель решает задачу классификации. Спектральный анализ отвечает на вопрос, как устроены её внутренние представления. Сопоставление этих двух типов данных позволяет аккуратнее интерпретировать архитектурные различия.


\section{Выводы по теоретической главе}

В данной работе теоретическая часть задаёт язык, на котором далее описываются методы и результаты. Понятия карты признаков, токена, патча, частотного радиуса, низких и высоких частот, спектрального сжатия и внутреннего представления используются в экспериментальной главе для построения единого протокола анализа.

Основная идея состоит в том, чтобы сравнить не только точность моделей, но и характер их внутренних представлений. Для этого ResNet, ViT и Swin рассматриваются как разные способы преобразования пространственного сигнала. ResNet представляет свёрточную иерархию, ViT — глобальную токенную обработку, Swin — локально-иерархический трансформерный подход. Спектральный анализ позволяет выразить различия между этими архитектурами через распределение энергии по пространственным частотам.

Таким образом, теоретическая база работы объединяет три направления: архитектуры компьютерного зрения, спектральный анализ и исследования spectral bias. Их сочетание позволяет сформулировать и проверить гипотезы о том, как архитектура влияет на спектральную динамику представлений по глубине.
